% !TEX root = ../my-thesis.tex
%
\chapter{Method: The Pipeline}
\label{sec:method}

This chapter describes the pipeline that turns raw debate text into a formally
evaluated Dung framework. Section~\ref{sec:method:overview} gives an overview,
and the remaining sections describe the five steps in turn, each with the design
decision behind it.

\section{Overview}
\label{sec:method:overview}

The pipeline runs in five steps. Step one acquires arguments and atomizes them
into self-contained claims. Step two examines every pair of claims and proposes
categorized attack candidates. Step three filters those candidates by confidence
and verifies the borderline cases. Step four converts the accepted candidates
into binary Dung attack facts. Step five solves the resulting framework under
grounded, preferred, and stable semantics with a Prolog solver.

\begin{figure}[t]
  \centering
  \begin{tikzpicture}[
    node distance = 6mm and 10mm,
    box/.style = {draw, rounded corners=3pt, fill=blue!8,
                  minimum width=82mm, minimum height=8mm,
                  align=center, font=\small},
    arr/.style = {-Stealth, thick},
  ]
    \node[box]                  (atom)   {Step 1: Argument Acquisition and Atomization\\
                                          \textit{LLM zero-shot generation or corpus atomization}};
    \node[box, below=of atom]   (batch)  {Step 2: Attack Detection and Classification\\
                                          \textit{propose labeled attack candidates per pair}};
    \node[box, below=of batch]  (filter) {Step 3: Confidence Filtering and Verification\\
                                          \textit{route by confidence score}};
    \node[box, below=of filter] (aaf)    {Step 4: AAF Conversion\\
                                          \textit{prune attacks to traced att/2 facts}};
    \node[box, below=of aaf]    (prolog) {Step 5: Solving under Dung Semantics\\
                                          \textit{grounded / preferred / stable extensions}};
    \draw[arr] (atom)   -- (batch);
    \draw[arr] (batch)  -- (filter);
    \draw[arr] (filter) -- (aaf);
    \draw[arr] (aaf)    -- (prolog);
  \end{tikzpicture}
  \caption{Overview of the pipeline, from raw debate text to formally evaluated
  Dung extensions in five steps.}
  \label{fig:method:pipeline}
\end{figure}

\section{Step 1: Argument Acquisition and Atomization}
\label{sec:method:atomize}

Arguments enter the pipeline in one of three ways. They are generated from
scratch by prompting an LLM to produce $n$ pro and $n$ con arguments on a topic,
extracted from an existing text corpus, or loaded directly from an argument
database. In all cases an argument is represented by a tuple of an identifier, a
stance label, and the argument text.

Natural-language arguments frequently bundle several claims. A compound sentence
such as ``Although nuclear energy produces waste, it is cleaner than coal''
mixes a concession with a main claim, and treating it as one unit would distort
the attack relations extracted next. The pipeline therefore prompts an LLM to
split each argument into self-contained atomic claims under three rules.

\begin{enumerate}
  \item \textbf{Drop concessive fragments.} Every component introduced by
    ``although'', ``while'', ``despite'', and similar markers is a rhetorical
    hedge rather than a standalone argument and is dropped.
  \item \textbf{Resolve pronouns.} A pronoun that clearly refers to a specific
    subject is replaced by that subject. A fragment whose referent cannot be
    recovered is removed, so that every atom is understandable in isolation.
  \item \textbf{Filter non-argumentative content.} A statement that neither
    supports nor opposes the claim under debate is dropped.
\end{enumerate}

A fourth, implicit rule keeps the stance intact. Reported or conceded positions
of the opposing side, such as ``opponents argue X'', are not promoted to their
own atom, because read in isolation they would carry the opposite stance of the
speaker. Each argument is atomized independently, which allows parallelization.
The resulting set of atomic arguments is denoted $\mathcal{A}'$ and forms the
input to the next step. Every atom keeps the identifier of its source argument
with a unique suffix, so that for instance \texttt{pro5} can yield the atoms
\texttt{pro5} and \texttt{pro5b}, which share the base \texttt{pro5}.

\section{Step 2: Attack Detection and Classification}
\label{sec:method:extraction}

With the atomic arguments in place, the pipeline examines every pair in
$\mathcal{A}'$ for a potential attack, regardless of stance, since same-stance
arguments can also conflict. Two pro-nuclear arguments may, for example,
prioritise incompatible safety strategies. Pairs of atoms that share a base are
skipped, since they originate from the same source argument.

Rather than a binary decision, the pipeline prompts the LLM to detect and
categorize each potential attack at once and to return a confidence score and a
justification. Asking for the type rather than a mere yes or no forces the model
to engage with the argumentative structure of the pair instead of pattern
matching on opposite stances, and the score and justification make each decision
traceable and provide a routing signal for the next step. Each candidate carries
four fields.

\begin{lstlisting}[caption={Fields of an attack candidate.},label={lst:method:candidate}]
direction:   A_ATTACKS_B | B_ATTACKS_A | NONE
attack_type: rebuttal | undercutting | undermining | NONE
confidence:  integer in [0, 100]
reason:      one-sentence justification
\end{lstlisting}

The prompt provides the full definitions of the three attack types from
Section~\ref{sec:foundations:typology} and enforces that if
\texttt{direction} is \texttt{NONE} then \texttt{attack\_type} is \texttt{NONE}
and \texttt{confidence} is zero, which keeps the output internally consistent.
Rebuttals are treated as symmetric, so the reverse edge is added automatically,
as motivated in Section~\ref{sec:foundations:typology}. All $\binom{|\mathcal{A}'|}{2}$
pairs are grouped into batches of size $b$, which reduces the number of API
calls from quadratic to $\binom{|\mathcal{A}'|}{2} / b$. Pairs for which the
model returns malformed output are retried individually.

This typed prompt is exactly the component whose value the ablation in
Section~\ref{sec:evaluation:ablation} measures, by comparing it against a binary
variant that omits the attack type.

\section{Step 3: Confidence Filtering and Verification}
\label{sec:method:filter}

Not all candidates are equally reliable. The confidence score, which is an
indication of the model's certainty rather than a calibrated probability, routes
each candidate through a three-way decision with two thresholds
$\theta_{\text{acc}}$ and $\theta_{\text{ver}}$.

\begin{equation}
\text{decision}(c) =
\begin{cases}
  \textbf{accept} & c \geq \theta_{\text{acc}} \\
  \textbf{verify} & \theta_{\text{ver}} \leq c < \theta_{\text{acc}} \\
  \textbf{reject} & c < \theta_{\text{ver}}
\end{cases}
\end{equation}

Borderline candidates in the verify region are re-examined by a dedicated
verifier prompt that asks the model to confirm or reject the attack given the two
arguments and the proposed type, but without access to the original confidence
score or justification. Withholding both keeps the verifier from anchoring on the
extractor's judgment. Accepted and confirmed candidates pass to the next step,
and rejected ones are discarded.

\section{Step 4: AAF Conversion}
\label{sec:method:conversion}

Every accepted candidate is converted to a binary Dung fact
$\texttt{att}(A, B)$, deliberately dropping the attack type from the formal
semantics. Unlike gradual QBAF semantics, which assign each argument a numerical
strength, Dung semantics ask which arguments can be jointly accepted, a set-based
rather than numerical question, so binary attack facts are the natural
representation. The attack type, the confidence score, and the justification are
not lost. They are retained as comments and metadata and remain available for
analysis and error inspection.

\section{Step 5: Solving under Dung Semantics}
\label{sec:method:solver}

The framework is solved by an AAF solver implemented in Prolog under the three
semantics of Section~\ref{sec:foundations:aaf}, using standard fixed-point and
enumeration procedures~\cite{cerutti2014}. The grounded extension is computed by
iterating the characteristic function from the empty set in polynomial time.
Computing preferred and stable extensions requires enumerating conflict-free
sets, which is exponential in the worst case.

To keep this tractable the solver applies a \emph{grounded-first reduction}. The
grounded extension and the arguments it attacks are fixed before the search, so
the exponential enumeration is restricted to the undecided arguments that remain.
This reduction is key to the scalability of the pipeline. On the frameworks in
the evaluation it brings the preferred and stable computation down to a fraction
of a second per topic, where a naive enumeration times out. The solver is
parametrised by the framework file and returns its results as JSON, which the
Python driver reads back.
